\section{Introduction}

Edge applications increasingly span devices that do not share a continuously connected network. Vehicle-mounted compute nodes, field tablets, and regional services may alternate between full connectivity, mesh-only connectivity, and total isolation. Public-safety mesh measurements show that such partitions can be frequent rather than exceptional in forested operational terrain \cite{zimbelman2022mesh}, while broader emergency-communications surveys describe heterogeneous LTE, ad hoc, and mission-critical network conditions \cite{kumbhar2015legacy,hoyhtya2018critical,wang2023overview}. A useful edge architecture must therefore specify not only how data moves when links are present, but also which writes remain meaningful when links are absent.

The dominant replication patterns make different trade-offs. Consensus-based state-machine replication, including Paxos- and Raft-style protocols, provides a single ordered log when a quorum is available \cite{lamport1998parttime,ongaro2014raft}. Multi-primary replication maximizes local write availability but pushes conflict resolution into reconciliation after partitions \cite{mucha2023database}. CRDT and local-first systems preserve availability for data types whose updates can be merged without coordination \cite{kleppmann2019localfirst,almeida2014delta}. These patterns are powerful, but they do not directly capture a class of edge applications in which events are authority-bearing: an event is not merely a state update, but an ordered decision made by a designated operational principal.

This paper uses incident-information management as a concrete instance of that problem. In command-hierarchical emergency response, some field observations can be queued locally, but authoritative incident state must be sequenced by the current command role. A wrong merge can be worse than a delayed commit if it produces inconsistent resource assignments, status transitions, or hazard annotations. The application domain motivates the design, but the systems question is broader: how should authority-bearing data be synchronized across intermittently connected edge tiers?

The contribution is a tiered synchronization architecture for authority-bearing data over intermittent multi-tier connectivity. The architecture combines a single-writer incident journal at the command vehicle, durable forward-only outboxes at responder vehicles and tablets, and explicit manual promotion rather than automatic leader election. It is structurally related to shared-log systems such as Tango \cite{balakrishnan2013tango}, but differs in why the writer exists: the writer is an operational authority selected by the incident workflow, not merely an elected replica leader. The submission version is being migrated toward an NCA evaluation that will add safety verification and baseline comparison results.

The rest of the paper is organized as follows. \Cref{sec:related} positions the design against state-machine replication, local-first systems, and edge-connectivity work. \Cref{sec:architecture} presents the three-tier architecture and fault model. \Cref{sec:sync} specifies the synchronization protocol. \Cref{sec:implementation} summarizes the prototype scaffold. \Cref{sec:validation} marks the planned NCA evaluation and remaining measurement work.
